\begin{table}[tbp]
\centering
\small
\caption{Standing representation notes for the closure formalization.}
\label{tab:representation-notes}
\begin{tabularx}{\textwidth}{@{}>{\raggedright\arraybackslash}p{0.25\textwidth}
>{\raggedright\arraybackslash}p{0.30\textwidth}
>{\raggedright\arraybackslash}X@{}}
\toprule
Representation issue & Where it appears & Disclosure effect \\
\midrule
Import structures & \(\chiCTp\), \(\AofE\), Beilinson, \(\eta\), and
cascade imports & Arithmetic theorem stacks are hypothesis fields in typed
records, not Lean axioms or internal derivations. \\
\addlinespace
Recognition carriers & \(\GammaPD\), \(\GammaSP\), \(\GammaGZ\) & The
arithmetic payload is supplied closure content and is consumed by the shell
as a hypothesis. \\
\addlinespace
Projection-packaged landing & Supporting conditionals and
\Cref{thm:closure:strong-bsd-conditional} & The formalization checks the
conditional projection/implication chain over supplied records, not an
independent arithmetic proof. \\
\addlinespace
Local \(\RefStable\) predicate & AOR-instance theorem & The Lean theorem
checks the paper-local syntactic predicate, not the full AOR meta-theory. \\
\addlinespace
Cross-paper \(\kapr\) import & \(\Cref{sec:aor-instance}\) & The
normalization is reused from the apparatus paper and cited there, not
re-mechanized in the closure axis. \\
\bottomrule
\end{tabularx}
\end{table}
